\section*{3. 对完全交的应用}

\subsection*{3.0.}
设 $X$ 是域 $k$ 上 $\mathbb P^N$ 中维数为 $n\geq1$ 的光滑完全交.
对 $\ell\ne\operatorname{char}(k)$, 弱 Lefschetz 定理及对偶性给出 (XI~1.6)
\begin{equation}\tag{3.0.1}
H^i(\overline X,\mathbb Q_\ell)=
\begin{cases}
0 & \text{若 $i$ 为奇数且 $i\ne n$},\\[1.5ex]
\mathbb Q_\ell\!\left(-\dfrac{i}{2}\right)
  & \text{若 $i$ 为偶数, $i\leq2n$ 且 $i\ne n$},\\[1.5ex]
H^n(\overline X,\mathbb Q_\ell) & \text{若 $i=n$}.
\end{cases}
\end{equation}
当 $k=\mathbb F_q$ 时, 这给出 zeta 函数的下述公式:
\begin{equation}\tag{3.0.2}
Z(t,X/\mathbb F_q)\prod_{i=0}^{n}(1-q^it)
=\det\!\left(1-t\varphi_q^{-1}\mid
\operatorname{Prim}^n(\overline X,\mathbb Q_\ell)\right)^{(-1)^{n+1}}.
\end{equation}
这里使用了本原分解及 (3.0.1).

\subsubsection*{推论 3.0.3.}
在 3.0 的假设下, 令 $k=\mathbb F_q$, 则有模 $p$ 同余
\[
\det\!\left(1-t\varphi_q^{-1}\mid
H^n(\overline X,\mathbb Q_\ell)\right)
\equiv
\det\!\left(1-t\mathfrak F_q\mid H^n(X,\mathcal O_X)\right).
\]

\medskip
\noindent\emph{证明.}
由同余公式 (XXII~3.2.1), 又由于
\begin{equation}\tag{3.0.4}
H^i(X,\mathcal O_X)\simeq
\begin{cases}
\mathbb F_q & i=0,\\[1.5ex]
0 & 1\leq i\leq n-1,\\[1.5ex]
0 & i\geq n+1,
\end{cases}
\end{equation}
有
\begin{equation}\tag{3.0.5}
Z(t,X/\mathbb F_q)\equiv
(1-t)^{-1}
\det\!\left(1-t\mathfrak F_q\mid H^n(X,\mathcal O_X)\right)^{(-1)^{n+1}}
\pmod p.
\end{equation}
因此由 (3.0.5) 和 (3.0.2) 得
\begin{equation}\tag{3.0.6}
\det\!\left(1-t\varphi_q^{-1}\mid
\operatorname{Prim}^n(\overline X,\mathbb Q_\ell)\right)
\equiv
\det\!\left(1-t\mathfrak F_q\mid H^n(X,\mathcal O_X)\right)
\pmod p.
\end{equation}
所求结论由下式得出
\begin{equation}\tag{3.0.7}
H^n(\overline X,\mathbb Q_\ell)\simeq
\begin{cases}
\operatorname{Prim}^n(\overline X,\mathbb Q_\ell)
  & \text{若 $n$ 为奇数},\\[2mm]
\operatorname{Prim}^n(\overline X,\mathbb Q_\ell)
 \oplus\mathbb Q_\ell\!\left(-\dfrac n2\right)
  & \text{若 $n$ 为偶数}.
\end{cases}
\end{equation}

\subsubsection*{推论 3.1.}
设 $X$ 是 $\mathbb F_q$ 上维数 $n\geq1$ 的光滑完全交. 若
\begin{equation}\tag{3.1.0}
H^n(\overline X,\mathbb Q_\ell)=N^1H^n(\overline X,\mathbb Q_\ell),
\end{equation}
则
\begin{equation}\tag{3.1.1}
\mathfrak F_p\quad\text{在下述空间上幂零:}\quad H^n(X,\mathcal O_X).
\end{equation}

\medskip
\noindent\emph{证明.}
由推论~2.2, (3.1.0) 蕴涵
\begin{equation}\tag{3.1.2}
\det\!\left(1-t\varphi_q^{-1}\mid
H^n(\overline X,\mathbb Q_\ell)\right)
\equiv1\pmod p
\qquad\text{(事实上, 甚至模 $q$ 成立)}.
\end{equation}
结合推论~3.0.3, 得到
\begin{equation}\tag{3.1.3}
\det\!\left(1-t\mathfrak F_q\mid H^n(X,\mathcal O_X)\right)=1.
\end{equation}
因此 $\mathfrak F_q$ 在 $H^n(X,\mathcal O_X)$ 上幂零, 从而 $\mathfrak F_p$ 也如此.
推论得证.

\medskip
\noindent\emph{底本注.}
照录法文层保留了印本的词序 \emph{intersection complète dans $\mathbb P^N$ lisse}、
误文 \emph{da décomposition primitive}, 以及 (3.0.4) 中
$1\leq{}\leq n-1$ 所漏的指标. 校订法文层和英文层采用通常词序及
\emph{la décomposition primitive}, 并在上同调范围中补回必需的指标 $i$.
